\documentclass{rtl}
\title{ROS2 \& Robotika Otonom dari Nol}
\subtitle{Silabus kursus berbasis simulasi --- cukup laptop biasa}
\doctype{Silabus Kursus}
\date{Versi 0.1 · November 2026}

\begin{document}
\makeheader

\begin{sorotan}[title={Ringkasan kursus}]
\begin{tabularx}{\linewidth}{@{}l X@{}}
\textbf{Untuk siapa} & Mahasiswa teknik, anggota tim robot, guru/pelajar SMK, dan pembelajar mandiri yang ingin masuk ke robotika otonom. \\
\textbf{Prasyarat} & Bisa memakai komputer dan pernah menulis program sederhana (bahasa apa saja). Linux dan Python diajarkan dari awal. \\
\textbf{Perangkat} & Laptop RAM 8\,GB (disarankan 16\,GB), Ubuntu 24.04 atau Windows + WSL2. Tanpa GPU khusus, tanpa robot fisik. \\
\textbf{Perangkat lunak} & ROS~2 Jazzy (LTS), Gazebo Harmonic, RViz2, Python (\texttt{rclpy}), OpenCV, Nav2, slam\_toolbox. Webots sebagai alternatif untuk laptop yang lebih lemah. \\
\textbf{Durasi} & 12 modul + proyek akhir, $\pm$\,40 jam belajar (video, latihan, kuis). \\
\textbf{Hasil} & Portofolio GitHub berisi paket ROS~2 dan robot otonom simulasi yang bisa memetakan, bernavigasi, dan mengenali objek. \\
\end{tabularx}
\end{sorotan}

\section*{Filosofi: simulasi dulu, hardware kemudian}
Robot sungguhan mahal, mudah rusak, dan jumlahnya terbatas --- akibatnya kebanyakan
mahasiswa hanya ``menonton'' satu-dua orang yang memegang robot. Di kursus ini setiap
peserta punya robotnya sendiri di dalam simulator. Konsep dan kode yang dipakai
(ROS~2, TF, Nav2, OpenCV) sama persis dengan yang dipakai di robot fisik, sehingga
mudah dipindahkan ke hardware nanti.

\section*{Tujuan pembelajaran umum}
Setelah menyelesaikan kursus, peserta mampu:
\begin{enumerate}
  \item Menyiapkan lingkungan kerja ROS~2 dan bekerja nyaman di terminal Linux.
  \item Merancang sistem robot sebagai kumpulan \emph{node} yang berkomunikasi lewat topic, service, dan action.
  \item Membuat paket ROS~2 dengan Python, lengkap dengan parameter dan file launch.
  \item Memodelkan robot (URDF/xacro), memahami kerangka koordinat (TF2), dan mensimulasikannya di Gazebo.
  \item Mengolah data sensor (LiDAR, kamera, IMU) dan menulis pengendali sederhana.
  \item Menerapkan computer vision dasar untuk mendeteksi objek dan penanda.
  \item Membangun peta (SLAM) dan menjalankan navigasi otonom dengan Nav2.
  \item Mendokumentasikan dan mempublikasikan proyek robotika di GitHub.
\end{enumerate}

\section*{Rincian modul}

\newcommand{\modul}[5]{%
  \subsection*{#1}
  {\small\color{rtlgray}Estimasi: #2}\par
  \textbf{Tujuan:} #3\par
  \textbf{Materi:} #4\par
  \textbf{Latihan simulasi:} #5\par}

\modul{Modul 0 --- Persiapan Lab di Laptopmu}{2 jam}
{Peserta memiliki lingkungan ROS~2 yang berjalan dan bisa memverifikasinya sendiri.}
{Pilihan instalasi (Ubuntu native, dual boot, WSL2 di Windows, Docker); instalasi ROS~2 Jazzy dan Gazebo Harmonic; tips laptop RAM 8\,GB; pemecahan masalah umum.}
{Menjalankan demo \texttt{talker}/\texttt{listener} dan membuka Gazebo dengan robot contoh. Lulus jika tangkapan layar ``cek kesehatan'' sudah hijau semua.}

\modul{Modul 1 --- Linux \& Python Secukupnya untuk Robotika}{3 jam}
{Peserta lancar memakai terminal dan menulis skrip Python yang rapi.}
{Navigasi berkas, izin, variabel lingkungan, \texttt{bash}; Git dasar; Python: fungsi, kelas, modul, \texttt{pip}/\texttt{venv}; membaca pesan error.}
{Skrip Python yang membaca berkas log sensor (CSV) dan menghitung rata-rata/maksimum; diunggah ke repo GitHub pribadi.}

\modul{Modul 2 --- Konsep Inti ROS~2: Node \& Topic}{3 jam}
{Peserta memahami arsitektur ROS~2 dan membuat publisher/subscriber sendiri.}
{Graph ROS~2, DDS secara singkat, node, topic, tipe pesan, QoS dasar; alat CLI (\texttt{ros2 topic}, \texttt{rqt\_graph}); turtlesim.}
{Mengendalikan turtlesim dengan node Python: menggambar persegi, lalu spiral; node kedua memantau posisi kura-kura.}

\modul{Modul 3 --- Service, Action, Parameter \& Launch}{3 jam}
{Peserta memilih pola komunikasi yang tepat dan menjalankan banyak node sekaligus.}
{Service vs action vs topic; parameter dan file YAML; file launch Python; namespace dan remapping.}
{Service ``reset arena'', action ``pergi ke titik'' dengan umpan balik progres, dan satu file launch yang menjalankan semuanya dengan parameter berbeda.}

\modul{Modul 4 --- Workspace \& Paket ROS~2 yang Rapi}{3 jam}
{Peserta membuat paket yang bisa dibangun, diuji, dan dipakai ulang orang lain.}
{Struktur workspace, \texttt{colcon build}, \texttt{package.xml}/\texttt{setup.py}, pesan kustom, pengujian sederhana, konvensi penamaan.}
{Paket \texttt{robot\_saya} dengan pesan kustom \texttt{StatusBaterai}, node simulator baterai, dan node peringatan.}

\modul{Modul 5 --- Model Robot: URDF, TF2 \& RViz}{4 jam}
{Peserta memodelkan robot beroda dan memahami hubungan antar kerangka koordinat.}
{URDF dan xacro, link dan joint, \texttt{robot\_state\_publisher}, TF2 (\texttt{map}--\texttt{odom}--\texttt{base\_link}), visualisasi di RViz2.}
{Merancang robot diferensial sendiri (ukuran, roda, sensor) dan memeriksa pohon TF-nya.}

\modul{Modul 6 --- Simulasi di Gazebo}{4 jam}
{Peserta menjalankan robotnya di dunia simulasi lengkap dengan sensor.}
{Dunia SDF, plugin diff-drive, sensor LiDAR, kamera, IMU; \texttt{ros\_gz\_bridge}; waktu simulasi; menjaga simulasi tetap ringan di laptop.}
{Membangun arena ``lorong kampus'', memasang LiDAR dan kamera, lalu menggerakkan robot dengan teleop keyboard.}

\modul{Modul 7 --- Membaca Sensor \& Kontrol Gerak}{4 jam}
{Peserta menulis perilaku otonom pertama berbasis data sensor.}
{Odometri dan sumber galatnya; pengolahan LaserScan; kontrol P/PID; mesin keadaan (state machine) sederhana.}
{Robot penghindar rintangan, lalu robot penyusur dinding (wall-following) dengan PID --- dinilai dari waktu tempuh dan jumlah tabrakan.}

\modul{Modul 8 --- Computer Vision untuk Robot}{4 jam}
{Peserta mengolah citra kamera robot untuk mengambil keputusan.}
{\texttt{cv\_bridge}, OpenCV: ruang warna HSV, thresholding, kontur; penanda ArUco dan estimasi pose; pengenalan singkat deteksi objek berbasis model pralatih (CPU).}
{Robot pengikut bola berwarna, lalu robot yang berhenti dan melapor saat melihat penanda ArUco tertentu.}

\modul{Modul 9 --- SLAM: Membuat Peta}{3 jam}
{Peserta membangun dan menyimpan peta lingkungan dari data LiDAR.}
{Konsep SLAM, occupancy grid, \texttt{slam\_toolbox}, menyimpan dan memuat peta, kualitas peta dan penyebab kegagalan.}
{Memetakan arena gudang simulasi dan membandingkan hasil dua strategi eksplorasi.}

\modul{Modul 10 --- Navigasi Otonom dengan Nav2}{4 jam}
{Peserta membuat robot bernavigasi sendiri ke tujuan di peta.}
{Lokalisasi (AMCL), costmap, planner dan controller, behavior tree, waypoint following; menyetel parameter dasar.}
{Robot mengunjungi 5 titik di gudang secara berurutan sambil menghindari rintangan yang muncul tiba-tiba.}

\modul{Modul 11 (Bonus) --- Sekilas Otonomi Drone dengan PX4}{3 jam}
{Peserta memahami bagaimana ROS~2 terhubung ke autopilot drone.}
{PX4 SITL dengan Gazebo, mode terbang, offboard control melalui ROS~2, keselamatan.\newline
\emph{Catatan: modul ini paling berat secara komputasi; tersedia rekaman demo bagi laptop yang tidak kuat.}}
{Lepas landas, terbang membentuk persegi, dan mendarat secara otonom di simulasi.}

\section*{Proyek akhir: Robot Kurir Kampus}
\begin{sorotan}[title={Brief proyek}]
Sebuah robot beroda harus mengantar ``paket'' di gedung kampus simulasi.
Robot memetakan gedung, lalu menerima daftar tujuan, bernavigasi ke tiap ruangan,
mengenali penanda ArUco di pintu untuk memastikan ruangan yang benar, dan
melaporkan status pengantaran ke sebuah topic. Rintangan dinamis (``orang lewat'')
muncul secara acak.
\end{sorotan}

\textbf{Yang dikumpulkan:} repo GitHub (paket ROS~2 + file launch), README dengan
cara menjalankan, video demo 2--3 menit, dan laporan singkat (1--2 halaman)
tentang keputusan desain dan kendala.

\begin{center}
\begin{tabularx}{0.92\linewidth}{@{}X r@{}}
\toprule
\textbf{Kriteria penilaian} & \textbf{Bobot} \\
\midrule
Fungsi: semua tujuan tercapai, robot tidak menabrak & 35\% \\
Persepsi: verifikasi ruangan dengan ArUco berjalan andal & 20\% \\
Kualitas kode: struktur paket, parameter, launch, keterbacaan & 20\% \\
Dokumentasi: README bisa diikuti orang lain dari nol & 15\% \\
Presentasi: video demo dan penjelasan keputusan desain & 10\% \\
\bottomrule
\end{tabularx}
\end{center}

Peserta yang lulus mendapat sertifikat penyelesaian, dan proyek terbaik ditampilkan
di repo dan kanal publik program (dengan izin peserta).

\section*{Format belajar}
\begin{multicols}{2}
\textbf{Kursus mandiri (evergreen)}
\begin{itemize}
  \item Video pendek 8--15 menit per topik
  \item Kode latihan dan solusi di GitHub
  \item Kuis singkat tiap modul
  \item Forum tanya jawab komunitas
\end{itemize}
\columnbreak
\textbf{Cohort live}
\begin{itemize}
  \item 5 minggu, 2 sesi live per minggu (malam hari WIB)
  \item Kelas kecil 20--30 peserta
  \item Review kode dan sesi tanya jawab
  \item Presentasi proyek akhir di minggu terakhir
\end{itemize}
\end{multicols}

\textbf{Pembagian cohort 5 minggu (usulan):}
Minggu 1: Modul 0--2 \,·\,
Minggu 2: Modul 3--5 \,·\,
Minggu 3: Modul 6--7 \,·\,
Minggu 4: Modul 8--10 \,·\,
Minggu 5: Proyek akhir + Modul 11 (bonus).

\end{document}
